% TOP_PROBLEM: {"id": 4700029, "title": "Around Kouchnirenko's conjecture II", "category_id": 4, "category": {"id": 4, "name": "algebra", "display_name": "Algebra", "description": "Group theory, ring theory, field theory, and algebraic structures.", "slug": "algebra", "order_index": 4, "created_at": "2026-07-31T15:26:25.671Z"}, "status": "partially_solved", "problem_number": "AMR-046-0029", "set_id": 13}
% FIRST_POSTED: 2026-10-02
% ENTRY_KIND: statement_only
% UnsolvedMath ID 4700029; source catalogue code AMR-046-0029
% !TeX program = lualatex
% Definition and sources only; this file is not an LLM solution attempt.
\documentclass[11pt,a4paper]{article}
\usepackage[margin=25mm]{geometry}
\usepackage{fontspec}
\setmainfont{lmroman10-regular.otf}[BoldFont=lmroman10-bold.otf,
 ItalicFont=lmroman10-italic.otf,BoldItalicFont=lmroman10-bolditalic.otf]
\usepackage{amsmath,amssymb,mathtools}
\usepackage{unicode-math}
\setmathfont{latinmodern-math.otf}
\AtBeginDocument{\let\setminus\smallsetminus}
\usepackage{xurl,hyperref}
\hypersetup{colorlinks=true,urlcolor=blue}
\setcounter{secnumdepth}{0}
\setlength{\parindent}{0pt}
\setlength{\parskip}{5pt}
\setlength{\emergencystretch}{3em}
\begin{document}
\section*{Around Kouchnirenko's conjecture II}\label{4700029}
{\small UnsolvedMath ID 4700029 · AMR-046-0029 · Algebra · AMR Open Problem Lists}
\subsection{Definitions and mathematical statement}
Let $m_1,m_2\ge1$, and let $f_i\in\mathbb R[x,y]$ have at most $m_i$
nonzero monomials with nonnegative integer exponents. Coefficients and
degrees are otherwise unrestricted. Count common zeros in the entire
real plane, including coordinate axes, but only when they are simple:
\[
 f_1(u,v)=f_2(u,v)=0,\qquad
 (f_1)_x(f_2)_y-(f_1)_y(f_2)_x\ne0
 \quad\text{at }(u,v).
\]
The proposed sharp bound is
\[
 \#\{\text{simple real common zeros}\}
       \le(2m_1-1)(2m_2-1).
\]
The question is whether this inequality holds for every allowed pair.
The count is of distinct isolated zeros, with no multiplicity assigned.
It is not limited to the positive quadrant, so positive-quadrant
counterexamples to the original Kouchnirenko bound do not automatically
settle this version.

The proposed value is attainable by separated-variable systems. For
example, choose distinct positive numbers $a_j$ and put
$g_m(t)=t\prod_{j=1}^{m-1}(t^2-a_j^2)$. This polynomial has $m$
nonzero monomials and $2m-1$ simple real roots. Taking
$f_1(x,y)=g_{m_1}(x)$ and $f_2(x,y)=g_{m_2}(y)$ attains the product.
Thus the unresolved direction is a universal upper bound for coupled
systems. A proof must include zeros on axes, where arguments involving
division by monomials require separate justification. The zero set may
contain other singular points or curves; none of those are counted, but
they are not excluded as additional features of an admissible system.
\subsection{Short English statement}
Is the product $(2m_1-1)(2m_2-1)$ the sharp number of simple real common zeros of two fewnomials?
\subsection{Sources}
\begin{itemize}
\item[S1] ULAM.AI and the UnsolvedMath Contributors, supplied problems.json, record 4700029 (AMR-046-0029), AMR Open Problem Lists. Catalogue identity and source transcription; CC BY 4.0.
\url{https://huggingface.co/datasets/ulamai/UnsolvedMath}.
\item[S2] Gasull - Some open problems in low dimensional dynamical systems (2020), Problem environment 29: Around Kouchnirenko's conjecture II. Original source indexed by AMR-046-0029.
\url{https://arxiv.org/abs/2012.02524}.
\end{itemize}
\end{document}
